\section{Conclusioni}
Ricapitolando, l'obbiettivo dell'esperienza è verificare la relazione che descrive l'energia di un fotone Compton in funzione dell'angolo di scattering e verificare la sezione d'urto espressa dall'equazione di Klein-Nishina. 

I risultati sperimentali ottenuti per il primo andamento e il test del $\chi^2$ dimostrano che la calibrazione e la disposizione angolare del set-up sperimentale sono state valutate correttamente e le misure effettuate sono coerenti con il modello teorico previsto. Tuttavia, per quanto riguarda la sezione d'urto sperimentale in funzione dei conteggi, il test del $\chi^2$ non è stato superato. 

Attraverso i fit parametrici si sono stimati i valori di $m_ec^2$ e del raggio dell'elettrone $r_e$ ottenendo: $m_ec^2=(523\pm12)\unit{\kilo\electronvolt}$ e $r_e=(3.01\pm0.65)\unit{fm}$. L'incertezza sul valore della massa a riposo dell'elettrone è chiaramente influenzato dalle larghezze gaussiane dei picchi misurati. La sua stima verrebbe dunque migliorata grazie ad un tempo di acquisizione maggiore, il quale aumenterebbe notevolmente la statistica. Si nota infatti che l'andamento gaussiano dei risulta maggiormente piccato all'aumentare della statistica acquisita.
La stima del raggio classico dell'elettrone risulta anch'essa confrontabile con il valore teorico ma caratterizzata da una incertezza piuttosto importante. Questo è infatti diretta conseguenza della misura sperimentale della sezione d'urto e le relative incertezze già precedentemente analizzate.

\vspace{0.5cm}
Infine, vengono qui riportate le possibili modifiche all'esperienza per migliorare la misura:
\begin{itemize}
    \item Lo studio manca di una analisi del fondo, potrebbe infatti aver influenzato il numero di conteggi osservati dal rivelatore.
    \item Come suggerito in sezione \ref{sez_err_angolo_solido} per ridurre l'errore sull'angolo solido coperto dal rivelatore sarebbe ragionevole proporre una geometria tale da aumentare le distanze coinvolte, a discapito di un rate di conteggi al secondo.
    \item Una ulteriore riflessione riguarda il bersaglio in quanto si potrebbe scegliere uno spessore particolarmente ridotto, rendendo così maggiormente confrontabili la lunghezza percorsa al suo interno dal fotone prima di interagire e lo spessore $dx$ presente nella sezione d'urto sperimentale.
\end{itemize}
